\SourceLeaf{318}%
for the sphere, which reduce to those same formulas when the eccentricity is set \(=0\) or the two semiaxes of the generating ellipse are set equal to each other and \(=1\).

As in Art. 3 above, I denote the latitude of the standard parallel on which the conformal mapping onto the sphere is based by \(P\) for the ellipsoid and by \(Q\) for the sphere. I also assume that this standard parallel is chosen so that \(Q\) equals the arithmetic mean of the latitudes of the two points concerned on the sphere. Let the latitudes themselves be \(Q+\frac12q\) and \(Q-\frac12q\). Let \(a,A,\alpha,e,\varphi,\theta\) retain the same meanings as in the first memoir, Arts. 2, 3, 4 and following. Here
\begin{list}{}{\setlength{\leftmargin}{1em}\setlength{\itemindent}{0pt}\setlength{\itemsep}{0pt}\setlength{\parsep}{0pt}}
\item \(a\) is the semimajor axis of the ellipsoid, or the radius of the equator,
\item \(A\) is the radius of the sphere,
\item \(1:\alpha\) is the constant ratio of longitude differences on the ellipsoid to the corresponding differences on the sphere,
\item \(e=\sin\varphi\) is the eccentricity of the generating ellipse,
\item \(\sin\theta=e\sin P\).
\end{list}

\indent I denote the great-circle arc between the two points on the sphere by \(s\), and the azimuths of this arc at the first and second endpoints by \(U+\frac12u\) and \(U-\frac12u\pm180^0\). Taking into account that the difference in longitude between the two points is \(=\alpha l\) gives first the four exact formulas
\[
\begin{aligned}
\sin\tfrac12s\mathord{.}\cos U&=\cos\tfrac12\alpha l\mathord{.}\sin\tfrac12q\\
\sin\tfrac12s\mathord{.}\sin U&=\sin\tfrac12\alpha l\mathord{.}\cos Q\\
\cos\tfrac12s\mathord{.}\cos\tfrac12u&=\cos\tfrac12\alpha l\mathord{.}\cos\tfrac12q\\
\cos\tfrac12s\mathord{.}\sin\tfrac12u&=\sin\tfrac12\alpha l\mathord{.}\sin Q
\end{aligned}
\]
and from these the approximate formulas
\[
\begin{aligned}
q&=s\mathord{.}\cos U(1+\tfrac1{24}qq-\tfrac1{24}ss+\tfrac18\alpha\alpha ll)\quad .\ .\ .\ .\ .\ .\ .\quad(1)\\
\alpha l&=s\mathord{.}\frac{\sin U}{\cos Q}(1-\tfrac1{24}ss+\tfrac1{24}\alpha\alpha ll)\quad .\ .\ .\ .\ .\ .\ .\ .\ .\ .\quad(2)\\
u&=s\mathord{.}\sin U\mathord{.}\operatorname{tang}Q(1+\tfrac1{12}ss+\tfrac1{24}uu)\quad .\ .\ .\ .\ .\ .\ .\ .\quad(3)
\end{aligned}
\]
It hardly needs pointing out that \(l,q,s,u\) in these three equations are understood to be expressed in units of the radius. It is easy to see that the equations are accurate
